\section{Columnar storage}
\label{sec:pg-storage-columnar}

ClickBench is a benchmark for analytical DBMS and it also reported results for both Hydra and Citus. Figure~\ref{fig:storage-clickbench-motivation} uses historical benchmark submissions on AWS machines~\autocite{clickbench_2023_columnar}. 
The benchmark consisted of around 43 queries over a dataset of around 100 million rows (3 runs with minum of second two taken as into the sum).

\begin{figure}[htbp]
\centering
\begin{tikzpicture}[x=1cm,y=1cm,font=\footnotesize]
  \node at (5.2,4.0) {Hot-query time (min)};
  \node at (11.2,4.0) {Stored size (GB)};
  \foreach \v in {0,40,80,120,160} {
    \draw[black!12] (3+.0275*\v,.1)--(3+.0275*\v,3.55);
    \node[below] at (3+.0275*\v,.05) {\v};
  }
  \foreach \v in {0,20,40,60,80} {
    \draw[black!12] (9+.055*\v,.1)--(9+.055*\v,3.55);
    \node[below] at (9+.055*\v,.05) {\v};
  }
  % From the pinned ClickBench result files: sum(min(run2,run3))/60;
  % storage bytes / 1e9. All 43 queries are included, without penalties.
  \foreach \y/\t/\d/\tl/\dl/\col in {
    3.1/157.0344/77.797067741/157.0/77.8/storageBlue,
    1.9/38.596661767/18.980918899/38.6/19.0/storageOrange,
    .7/7.983093133/18.995669982/8.0/19.0/storageGreen} {
    \draw[draw=\col,fill=\col!25]
      (3,\y-.25) rectangle (3+.0275*\t,\y+.25);
    \node[anchor=west] at (3+.0275*\t+.08,\y) {\tl};
    \draw[draw=\col,fill=\col!25]
      (9,\y-.25) rectangle (9+.055*\d,\y+.25);
    \node[anchor=west] at (9+.055*\d+.08,\y) {\dl};
  }
  \node[anchor=east,align=right] at (2.75,3.1) {PostgreSQL 14\\2022 submission};
  \node[anchor=east,align=right] at (2.75,1.9) {Citus 11.0\\2022 submission};
  \node[anchor=east,align=right] at (2.75,.7) {Hydra Columnar\\2023 submission};
  \draw[black!55] (3,.1)--(7.4,.1);
  \draw[black!55] (9,.1)--(13.4,.1);
  \node at (5.2,-.65) {Sum over all 43 queries};
  \node at (11.2,-.65) {$1\ \mathrm{GB}=10^9$ bytes};
\end{tikzpicture}
\caption{Historical ClickBench, calculated from the
repository snapshot of 20 September 2023. Each submission uses
\texttt{c6a.4xlarge} machine with 16 vCPUs, 32~GiB RAM, and 500~GB of GP2
storage. PostgreSQL and Citus
results are dated 1 July 2022, and Hydra results 19 September 2023.
The times use the sum of shorter of runs two and three for every query. Note that these are "best-fit" results so the configuration was different. It's good however for quick comparison of performance.}
\label{fig:storage-clickbench-motivation}
\end{figure}

Both columnar submissions store about 19~GB, compared with PostgreSQL's
77.8~GB. This was due to support of compression on stored columns. Hydra~\autocite{hydra_columnar_source} had much better performance due to introduction of vectorized SQL execution for aggregates and column cache that wasn't present in Citus~\autocite{citus_distributed_2021}.
These results do not translate directly to BioseqDB as the benchmarks tests some arbitrarily-built queries, however it motivated me to put big effort into modernizing Hydra codebase and incorporating it into BioseqDB as mentioned in previous section~\ref{sec:pg-storage-role}.

\subsection{Columnar layout}


% %Stripe size is not configurable and is set to size of disk block minus page header so it can exactly fit into one disk block (around $8KB$, but depends on compile-time value Postgres version provides).
% \item A \textbf{stripe} is a batch of appended values, where values can be anything. In case of the user table these are grouped metadata values (native Postgres types) and for sequence columns - sequence locators. For internal sharded storage tables these are sequence fragments itself.
%     It's important to note that TOAST can have a limit up to $1GB$ per row, so large values (not only long sequence but also if user stores large BLOB) can span multiple stripes.
  

The columnar layout has three levels:

\begin{itemize}
  \item A \textbf{stripe} is a batch of appended rows. The default maximum
    is 150,000 rows.
  \item A \textbf{chunk group} is a smaller row interval within a stripe.
    The default maximum is 10,000 rows. Every column uses the same group
    boundaries.
  \item A \textbf{column chunk} contains one column's values for one chunk
    group. It has a presence bitmap and a value buffer. A presence bit says
    whether that row contains a non-NULL value.
\end{itemize}

Stripe and chunk-group sizes are table options. Smaller groups give finer
read selection, while larger groups give the compressor more values to work
with. This is the grouping model described by Citus
Columnar~\autocite{citus_columnar_compression_2021}.

Figure~\ref{fig:storage-heap-to-stripe} shows the same six rows in both
layouts. A heap tuple stores \texttt{id}, \texttt{sample\_id}, and the
sequence locator together. The stripe keeps all three columns of those rows.
Each column occupies its own sequence of chunks. A chunk group is only a
shared row interval: group 0 holds rows 1--3 of every column, and group 1
holds rows 4--6. The drawing uses three rows per group. The table defaults
are 10,000 rows per group and 150,000 rows per stripe. The last group of a
stripe may be shorter if fewer rows remain.

The writer fills one in-memory group for every column, serializes each
column chunk when the group is full, and flushes the stripe when the stripe
row limit is reached. Mapping those concatenated buffers onto PostgreSQL
pages is a later step in Section~\ref{sec:pg-storage-pages}. A stripe is
therefore not one disk page.

The callout in the figure is the stored form of one chunk, here \texttt{id}
of group 0. The presence bitmap has one bit per row in the group. The value
buffer holds only the present values. The skip record stores the smallest
and largest non-NULL values and the row count, which later pruning uses.
Offsets of the two buffers are also stored so a reader can open this chunk
without decoding the previous one. A zero presence bit would omit that row
from the value buffer. Figure~\ref{fig:storage-column-streams} shows that
case with explicit NULL slots.

\begin{figure}[H]
\centering
\begin{tikzpicture}[x=1cm,y=1cm,font=\footnotesize]
  \node[anchor=west,font=\bfseries] at (0,6.10)
    {Heap tuples become a stripe of column chunks};
  \node[anchor=west,text width=4.4cm,align=left] at (0,5.55)
    {Heap: one tuple stores every column};
  \foreach \y/\id/\sid/\rn in {
    4.45/1/12/1,
    3.80/2/18/2,
    3.15/3/7/3,
    2.50/4/41/4,
    1.85/5/9/5,
    1.20/6/22/6} {
    \node[anchor=east] at (1.05,\y) {row \rn};
    \node[storage cell,minimum width=.62cm,fill=storageBlue!15]
      at (1.55,\y) {\id};
    \node[storage cell,minimum width=.62cm,fill=storageOrange!15]
      at (2.28,\y) {\sid};
    \node[storage cell,minimum width=.78cm,fill=storageGreen!15]
      at (3.15,\y) {loc};
  }
  \draw[storage arrow] (3.90,2.85)--(4.70,2.85);
  \node[align=center] at (4.30,3.35) {same\\rows};
  \draw[draw=black!45] (4.85,0.42) rectangle (12.35,5.48);
  \node[anchor=north west,font=\bfseries] at (5.05,5.20)
    {One stripe: every column of these rows};
  \node at (7.55,4.55) {\texttt{id}};
  \node at (9.65,4.55) {\texttt{sample\_id}};
  \node at (11.40,4.55) {\texttt{seq}};
  \node[anchor=west] at (5.05,3.25) {group 0};
  \node[anchor=west] at (5.05,1.175) {group 1};
  \node[storage cell,minimum width=1.70cm,minimum height=.95cm,
    fill=storageBlue!15,draw=storageBlue,line width=1.1pt]
    (idchunk) at (7.55,3.25) {1, 2, 3};
  \node[storage cell,minimum width=1.70cm,minimum height=.95cm,
    fill=storageOrange!15]
    at (9.65,3.25) {12, 18, 7};
  \node[storage cell,minimum width=1.70cm,minimum height=.95cm,
    fill=storageGreen!15]
    at (11.40,3.25) {loc, loc, loc};
  \node[storage cell,minimum width=1.70cm,minimum height=.95cm,
    fill=storageBlue!15]
    at (7.55,1.175) {4, 5, 6};
  \node[storage cell,minimum width=1.70cm,minimum height=.95cm,
    fill=storageOrange!15]
    at (9.65,1.175) {41, 9, 22};
  \node[storage cell,minimum width=1.70cm,minimum height=.95cm,
    fill=storageGreen!15]
    at (11.40,1.175) {loc, loc, loc};
  \node[draw=storageBlue,fill=storageBlue!6,minimum width=12.20cm,
    minimum height=1.48cm,inner sep=0pt,anchor=north west]
    (chunkdetail) at (0.15,0.08) {};
  \node[anchor=west,font=\bfseries,text=storageBlue] at (0.35,-0.12)
    {One column chunk};
  \node[anchor=west,text=black!65] at (4.55,-0.12)
    {\texttt{id} of group 0, stored};
  \node[anchor=east] at (2.15,-0.72) {presence bits};
  \foreach \i/\b in {0/1,1/1,2/1}
    \node[storage cell,minimum width=.55cm,fill=storageBlue!12]
      at (2.65+0.70*\i,-0.72) {\b};
  \node[anchor=east] at (6.35,-0.72) {value buffer};
  \foreach \i/\v in {0/1,1/2,2/3}
    \node[storage cell,minimum width=.55cm,fill=storageGreen!13]
      at (6.85+0.70*\i,-0.72) {\v};
  \node[anchor=east] at (9.55,-0.72) {aggregates};
  \node[storage cell,minimum width=2.15cm,fill=storageOrange!12]
    at (10.80,-0.72) {min $1$, max $3$};
  \draw[storage arrow,storageBlue]
    (idchunk.south) -- (7.55,2.20) -- (12.80,2.20) -- (12.80,-0.66)
    -- (chunkdetail.east);
  \node[fill=white,inner sep=2pt,text=storageBlue] at (10.05,2.20)
    {One column chunk};
\end{tikzpicture}
\caption{The same six rows as heap tuples and as one stripe. Matching
colours mark the same column. The stripe holds every column. The callout
is the stored \texttt{id} chunk of group 0: presence bits, present values,
and aggregates. Logical layout is shown, not the later byte order on disk.}
\label{fig:storage-heap-to-stripe}
\end{figure}

\subsection{Two independent ways to avoid reads}

\emph{Projection} selects the columns required by the query. The sample
query needs \texttt{id} for its output and \texttt{sample\_id} for its
condition. It does not need \texttt{seq}.

\emph{Min/max pruning} selects chunk groups. For an ordinary ordered column,
the writer records the smallest and largest non-NULL values in each chunk.
A condition can reject a group if none of its values could satisfy the
condition. For example, a group with sample identifiers from 1 to 100 cannot
satisfy \texttt{BETWEEN 500 AND 600}.

Figure~\ref{fig:column-store-layout} shows both decisions. A retained group
still needs filtering. Each retained row must be tested against the condition:
bounds of 450 and 700 do not say which rows fall between 500 and 600.

\begin{figure}[H]
\centering
\begin{tikzpicture}[x=1cm,y=1cm,font=\footnotesize]
  \node[anchor=west,font=\bfseries] at (0,4.45)
    {One stripe: three chunk groups and three columns};
  \node at (3.3,3.85) {\texttt{id}};
  \node at (6.15,3.85) {\texttt{sample\_id}};
  \node at (9.0,3.85) {\texttt{seq}};
  \foreach \y/\g in {3.05/0,1.95/1,.85/2} {
    \node[anchor=east] at (1.7,\y) {group \g};
    \foreach \x in {1.9,4.75,7.6}
      \draw[draw=black!45,fill=black!5] (\x,\y-.4)
        rectangle +(2.65,.8);
  }
  \draw[storageBlue,fill=storageBlue!15] (1.9,1.55) rectangle (4.55,2.35);
  \draw[storageBlue,fill=storageBlue!15] (4.75,1.55) rectangle (7.4,2.35);
  \node at (3.23,1.95) {read identifiers};
  \node at (6.08,1.95) {bounds $[450,700]$};
  \node at (6.08,3.05) {bounds $[1,100]$};
  \node at (6.08,.85) {bounds $[800,950]$};
  \foreach \y in {3.05,.85} {
    \draw[storageOrange,line width=1.1pt] (2.05,\y-.28)--(4.4,\y+.28);
    \draw[storageOrange,line width=1.1pt] (2.05,\y+.28)--(4.4,\y-.28);
    \node[anchor=west,text=storageOrange] at (10.55,\y) {skip group};
  }
  \node[align=center,text=black!65] at (8.93,1.95) {locator stream\\not requested};
  \node[anchor=west,text=storageBlue,align=left] at (10.55,1.95)
    {read group;\\test its rows};
  \draw[decorate,decoration={brace,mirror,amplitude=4pt},black!60]
    (7.6,.22)--(10.25,.22);
  \node[align=center] at (8.93,-.37) {no sequence\\fragment reads};
\end{tikzpicture}
\caption{Column selection and group pruning for the metadata query.
Blue cells are read. Orange crosses mark groups whose scalar bounds exclude
every row. Unrequested sequence locators remain unread. The drawing shows
logical chunks, not their byte order on disk.}
\label{fig:column-store-layout}
\end{figure}

Sequence columns are excluded from this min/max test. A locator's byte order
does not describe sequence similarity. A separate scalar length, date, or
sample identifier can supply useful bounds. When a supported comparison or
usable bound is missing, the reader keeps the group.

\subsection{A column chunk separates missing values from values}

Figure~\ref{fig:storage-column-streams} uses six rows to show the two
buffers in one column chunk. The presence bitmap has one bit per row.
Only present values enter the value buffer, with the byte alignment required by
their PostgreSQL types. The reader uses the bitmap to put values back into
their row positions.

\begin{figure}[H]
\centering
\begin{tikzpicture}[x=1cm,y=1cm,font=\footnotesize]
  \node[anchor=east] at (2.1,3.15) {row position};
  \node[anchor=east] at (2.1,2.35) {column value};
  \node[anchor=east] at (2.1,1.55) {presence bit};
  \foreach \i/\v/\b in {0/12/1,1/NULL/0,2/17/1,3/17/1,4/NULL/0,5/21/1} {
    \node at (2.75+1.35*\i,3.15) {\i};
    \node[storage cell,minimum width=1.2cm,fill=black!4]
      at (2.75+1.35*\i,2.35) {\v};
    \node[storage cell,minimum width=1.2cm,fill=storageBlue!12]
      at (2.75+1.35*\i,1.55) {\b};
  }
  \node[anchor=east] at (2.1,.25) {value buffer};
  \foreach \i/\v in {0/12,1/17,2/17,3/21}
    \node[storage cell,minimum width=1.2cm,fill=storageGreen!13]
      at (2.75+1.35*\i,.25) {\v};
  \draw[storage arrow] (8.0,.25)--(9.0,.25);
  \node[storage cell,fill=storageOrange!12,minimum width=2.25cm]
    at (10.25,.25) {compress values};
  \node[anchor=west,text=black!65] at (2.15,-.5)
    {Only four values are stored; the bitmap keeps their six row positions.};
\end{tikzpicture}
\caption{Presence and values have different jobs. Generic column compression
acts on the serialized value buffer. Sequence rank packing, introduced
earlier, is a separate encoding step.}
\label{fig:storage-column-streams}
\end{figure}

The value buffer can use PostgreSQL PGLZ, LZ4, or Zstandard, depending on the
build and table options. Metadata records the compression method used and the
decoded length. Successful LZ4 and Zstandard output is not always smaller
than its input, so compression does not guarantee a space saving for every
chunk. No particular compression ratio is assumed here.

Abadi et al. show that the choice of compression method depends on both
the data and the query workload~\autocite{abadi2006compression}.
For compressed columns, BioSeqDB reconstructs SQL values from decoded
buffers, obtained by decompression or cache reuse, so both stored size
and decoding work matter.

The physical order differs from the logical grid. Within a stripe, one
column's presence buffers are written first, followed by that column's value
buffers. The next column follows. Stored offsets locate each buffer without
requiring the reader to decode earlier columns. A column added after a stripe
was written can be supplied from its constant PostgreSQL default or as NULL.

\section{Store byte streams in PostgreSQL relation pages}
\label{sec:pg-storage-pages}

A compressed buffer need not fit in one database page. BioSeqDB gives column
writers a continuous logical byte space and maps it onto page payloads.
Block 0 contains the metapage, block 1 is reserved, and stream data begin in
block 2. Each data page keeps its PostgreSQL page header.

Let $\mathrm{BLCKSZ}$ be the PostgreSQL page size in bytes, $H$ the
page-header size, and $P=\mathrm{BLCKSZ}-H$ the remaining payload capacity.
Measure $x\geq0$ from the first stream byte in block 2.
Then
\[
  \operatorname{block}(x)=2+\left\lfloor x/P\right\rfloor,
  \qquad
  \operatorname{position}(x)=H+(x\bmod P).
\]
This choice of origin makes the stream's first byte $x=0$. Page headers do
not consume positions in the logical byte stream. Figure~\ref{fig:storage-pages}
shows a request that crosses a page boundary.

\begin{figure}[htbp]
\centering
\begin{tikzpicture}[x=1cm,y=1cm,font=\footnotesize]
  \node[anchor=west,font=\bfseries] at (0,4.05)
    {Toy pages: 2 header bytes and 10 payload bytes};
  \node[anchor=east] at (1.45,2.85) {logical bytes};
  \draw[draw=black!50,fill=black!5] (1.75,2.5) rectangle (11.75,3.2);
  \draw[storageBlue,fill=storageBlue!20] (5.75,2.5) rectangle (9.25,3.2);
  \foreach \x/\v in {1.75/0,5.75/8,6.75/10,9.25/15,11.75/20}
    \node[above=2pt] at (\x,3.2) {$\v$};
  \node at (7.5,2.85) {request $[8,15)$};
  \node[anchor=east] at (1.45,1.15) {stored pages};
  \draw[fill=black!20,draw=black!60] (1.75,.8) rectangle (2.55,1.5);
  \draw[fill=black!5,draw=black!60] (2.55,.8) rectangle (6.55,1.5);
  \draw[storageBlue,fill=storageBlue!20] (5.75,.8) rectangle (6.55,1.5);
  \draw[fill=black!20,draw=black!60] (7.0,.8) rectangle (7.8,1.5);
  \draw[fill=black!5,draw=black!60] (7.8,.8) rectangle (11.8,1.5);
  \draw[storageBlue,fill=storageBlue!20] (7.8,.8) rectangle (9.8,1.5);
  \node at (2.15,1.15) {H};
  \node at (7.4,1.15) {H};
  \node at (6.15,1.15) {2};
  \node at (8.8,1.15) {5 bytes};
  \node at (4.15,.35) {block 2};
  \node at (9.4,.35) {block 3};
  \draw[storage arrow,storageBlue] (6.25,2.45)--(6.15,1.55);
  \draw[storage arrow,storageBlue] (8.0,2.45)--(8.8,1.55);
  \node[anchor=west,text=black!65] at (1.75,-.3)
    {Copy two bytes, skip the next page header, then copy five bytes.};
\end{tikzpicture}
\caption{Logical offsets hide page boundaries. The example uses deliberately
small pages to make the split visible. Actual page sizes and header sizes
come from the PostgreSQL build.}
\label{fig:storage-pages}
\end{figure}

This mapping follows the Citus column-store design. The public source gives
the corresponding page and logical-offset operations.\coderef
{citusstorage/src/backend/columnar/columnar_storage.c}
{Citus columnar page mapping}{columnar_storage.c}{118}
PostgreSQL documents the common page header separately from the layout an
access method puts after it~\autocite{postgresql17_storage_page_layout}.

\subsection{Reserve space, write bytes, publish metadata}

Concurrent writers must not allocate the same row numbers or byte interval.
The metapage records the next stripe number, row number, and free byte offset.
The write path separates allocation from visibility. It has three
responsibilities:

\begin{itemize}
  \item \textbf{Reserve.} Under the relation extension lock, reserve
    non-overlapping identifiers and byte space.
  \item \textbf{Write.} Copy the stream bytes into relation pages. Page
    changes use PostgreSQL buffer locks and generic write-ahead logging
    (WAL).
  \item \textbf{Record completed data.} Store rows, offsets, lengths, and
    chunk metadata in the catalogs. The bytes and catalog changes belong
    to the same transaction. Other transactions can read a completed stripe
    when its committed metadata is visible under their snapshots.
\end{itemize}

A reserved byte range alone is not a readable stripe. If a write does not
complete, the range does not authorize readers to interpret unfinished
bytes as table rows. Generic WAL supplies page-change recovery for logged
relations~\autocite{postgresql17_generic_wal}; the storage code still has to
publish the right catalog state and release its locks correctly.

\subsection{Why use relations instead of private payload files?}

The primary column and fragment data use PostgreSQL relations. A read obtains
a shared database buffer. A \textbf{page pin} keeps that buffer assigned to
the page while it is in use; a buffer lock protects its contents during the
copy. After copying the requested bytes, the reader releases the pin and
lock. The copied column buffer belongs to the reader. A sequence view
therefore does not keep every database page it has touched pinned for the
rest of the query.

This choice reuses PostgreSQL's shared buffer manager, relation identity,
locking, and WAL path. A private-file implementation would need separate
rules connecting its writes, deletion, and recovery to database transactions.
These benefits do not establish every transaction property of a custom
access method automatically. They identify the machinery used for these
particular bytes. Derived index images can use a different read mechanism,
described in the next chapter.
